% 图51-F51 Q3 求解流程图 — 修复溢出版
\documentclass[border=5pt]{article}
\usepackage[margin=0pt,paperwidth=14cm,paperheight=16cm]{geometry}
\pagestyle{empty}
\usepackage{fontspec}
\usepackage{amsmath}
\usepackage{ctex}
\newcommand{\fontdir}{../../../00_知识库/论文写作/LaTeX模板_2026国赛版/fonts/}
\setCJKfamilyfont{song}[
  Path = \fontdir, UprightFont = SourceHanSerifCN-Regular.otf,
  BoldFont = SourceHanSerifCN-Bold.otf, AutoFakeBold = true,
]{SourceHanSerifCN}
\newcommand*{\song}{\CJKfamily{song}}
\usepackage{tikz}
\usetikzlibrary{arrows.meta,positioning,calc,shapes.geometric}

\definecolor{primary}{HTML}{1F3A93}
\definecolor{accent}{HTML}{D35400}
\definecolor{green}{HTML}{27AE60}
\definecolor{gray}{HTML}{95A5A6}
\definecolor{textDark}{HTML}{333333}

\tikzset{
  arr/.style = {-{Stealth[length=2.0mm,width=1.8mm]},line width=0.8pt,draw=textDark!60},
  arrD/.style = {-{Stealth[length=1.8mm,width=1.6mm]},line width=0.7pt,draw=accent,dashed,dash pattern=on 4pt off 2pt},
  inNode/.style = {trapezium,trapezium left angle=65,trapezium right angle=65,
    draw=primary,line width=0.9pt,fill=primary!10,text=textDark,
    minimum width=4.0cm,minimum height=0.8cm,align=center,
    font=\song\footnotesize,trapezium stretches=true},
  procNode/.style = {rectangle,draw=green,line width=0.8pt,fill=green!10,text=textDark,
    minimum width=3.4cm,minimum height=0.75cm,align=center,
    font=\song\footnotesize,rounded corners=2pt},
  decNode/.style = {diamond,draw=accent,line width=0.8pt,fill=accent!8,text=textDark,
    minimum width=2.8cm,minimum height=1.0cm,align=center,
    font=\song\footnotesize,aspect=1.8,inner sep=1pt},
  outNode/.style = {rectangle,draw=primary,line width=1.1pt,fill=primary!8,text=textDark,
    minimum width=3.6cm,minimum height=0.8cm,align=center,
    font=\song\footnotesize\bfseries,rounded corners=2pt},
}

\begin{document}\song
\begin{tikzpicture}[scale=0.82,node distance=1.2cm]

% ── 1. 输入 ──
\node[inNode] (Input) {输入：初值+边界（附件1）+尺寸（附件2）};

% ── 2. 空间离散 ──
\node[procNode,below=of Input] (Disc) {空间离散（$N=80$，元体平衡二阶）};

% ── 3. 单步隐式推进 ──
\node[procNode,below=of Disc] (Step) {单步隐式推进（追赶法求解）};

% ── 4. 判断 ──
\node[decNode,below=1.0cm of Step] (Check) {中心含水率\\[-1pt]$C(0) < 0.15$？};

% ── 5. 步进（否分支，右侧） ──
\node[procNode,right=2.8cm of Check] (Next) {步进下一输出步\\[-1pt]（自适应/固定步）};

% ── 6. 回退（是分支，下方） ──
\node[procNode,below=1.0cm of Check] (Back) {回退至该输出步起点\\[-1pt]保留子步历史};

% ── 7. 二分 ──
\node[procNode,below=of Back] (Bisect) {子步内二分回溯\\[-1pt]（精度 $10^{-3}$ s）};

% ── 8. 输出 ──
\node[outNode,below=of Bisect] (Output) {输出 $t_{\rm end}$ 写 result3.xlsx};

% ====== 主线 ======
\draw[arr] (Input) -- (Disc);
\draw[arr] (Disc) -- (Step);
\draw[arr] (Step) -- (Check);

% ====== 否分支：Check→Next→回流至Step右侧 ======
\draw[arr] (Check.east) -- node[above,font=\song\scriptsize] {否} (Next.west);
% 先竖直上升到Step.east同一高度，再水平向左接Step.east
\draw[arr,draw=gray!60!black,rounded corners=4pt] (Next.north)
  |- (Step.east)
  node[pos=0.75,above,font=\song\scriptsize,text=gray!70!black]
  {循环回路\\[-1pt]未达判据};

% ====== 是分支：Check→Back→Bisect→Output ======
\draw[arr] (Check.south) -- node[right,font=\song\scriptsize] {是} (Back.north);
\draw[arr] (Back) -- (Bisect);
\draw[arr] (Bisect) -- (Output);

% ====== 左侧二分回溯长虚线 ======
% 从Bisect左侧引出，向左走x=-3.3，向上折返到Step左侧
\draw[arrD] (Bisect.west) -- ++(-2.0,0)
  node[midway,above,font=\song\scriptsize,text=accent] {二分回溯}
  |- ([xshift=-0.3cm]Step.west);

\end{tikzpicture}
\end{document}

% ── 修订版改动清单 ──
% 1. paperwidth 从 16cm 缩到 14cm，paperheight 从 14cm 缩到 16cm（高一点防底部切）
% 2. scale 从 0.85 微调到 0.82，整体更紧凑
% 3. 节点 vertical node distance 从 1.6cm 压到 1.2cm，菱形/回退额外间距调到 1.0cm
% 4. 右侧 Next 节点 horizontal 距离从 4cm 缩到 2.8cm，回流箭头不再溢出
% 5. 左侧二分回溯虚线：x 步进从 -3.0 缩到 -2.0，始终在画布内
% 6. 所有节点 minimum width/height 缩小（inNode 4.0→4.0, procNode 3.8→3.4, decNode 宽 3.5→2.8, outNode 4.0→3.6）
% 7. arr 箭头线宽从 0.9pt 缩到 0.8pt，arrD 从 0.8pt 缩到 0.7pt
% 8. 右侧回流改为从Next左上角→左弯折→Step顶部，圆角弯折(rounded corners=3pt)，删除灰色虚线框
% 9. "循环回路未达判据"小灰字直接放在箭头右侧，无需外框